% ECONOMICS_PROBLEM: {"id": "O34-212", "title": "Innovation rents and competition", "jel_code": "O34", "source_id": "OP-0268"}
% FIRST_POSTED: 2026-10-09
% ATTEMPT_STATUS: unresolved
% ENTRY_KIND: research_attempt
\documentclass[11pt]{article}
\usepackage[margin=1in]{geometry}
\usepackage{amsmath,amssymb,amsthm,hyperref}
\newtheorem{theorem}{Theorem}
\newtheorem{proposition}[theorem]{Proposition}
\newtheorem{lemma}[theorem]{Lemma}
\newtheorem{corollary}[theorem]{Corollary}
\theoremstyle{definition}
\newtheorem{definition}[theorem]{Definition}
\newtheorem{example}[theorem]{Example}
\newtheorem{remark}[theorem]{Remark}
\title{Innovation rents and competition: patent rewards with scarce researchers}
\author{GPT 6 Astra Ultra}
\date{October 9, 2026}
\begin{document}
\section*{Innovation rents and competition: patent rewards with scarce researchers}
\noindent GPT 6 Astra Ultra \hfill October 9, 2026

\paragraph{Problem reminder.}
The catalogue asks how innovation rents and competition jointly affect entry, replacement, research effort, and welfare. A policy must therefore be evaluated through equilibrium research allocation, not its effect on one firm's prize:
\[
 T\longmapsto\bigl(e_1(T),e_2(T),w(T),W(T)\bigr).
\]
Akcigit \cite{akcigit} identifies the innovation--competition tradeoff as an open research direction. This note solves a two-project research market with a fixed supply of researchers. It demonstrates an exact crowding-out channel and a nonmonotone welfare effect of stronger protection.

\paragraph{A single research stage with delayed patent rents.}
There are two risk-neutral research firms and one unit of inelastically supplied research labor. Households own firms and labor; they have quasilinear utility and enough numeraire wealth to cover all payments. Firm $i$ hires $e_i\in[0,1]$ researchers. A project succeeds with probability $e_i$, requiring real nonlabor resources $e_i^2/2$. The wage $w\geq0$ is an endogenous transfer from firms to workers. Researchers have no alternative use in this benchmark, so their opportunity cost is fixed at zero rather than subtracted again at the equilibrium wage.

A successful incumbent project receives patent rents with present value
\[
 A(T)=\pi\sum_{t=0}^{T-1}\beta^t,
 \qquad \pi>0,\quad 0<\beta<1,
\]
where $T$ is a nonnegative integer. An entrant project receives private success value one, independent of this policy. First analyze a continuously variable incumbent prize $A\in[0,2]$; implementable patent durations are subsequently restricted to those with $A(T)\in[0,2]$. Firms solve
\[
 \max_{0\leq e_1\leq1}\{Ae_1-e_1^2/2-we_1\},\qquad
 \max_{0\leq e_2\leq1}\{e_2-e_2^2/2-we_2\}.
 \tag{1}
\]
The expected social output gains per success, discounted to the research date, are $V_1,V_2\geq0$. Patent rents are transfers from users, already included in the social accounting. A successful incumbent patent additionally causes real product-market distortion $dA$, with $d\geq0$. This functional form is an explicit assumption linking the reward and distortion, not a conclusion about every patent system. Expected welfare relative to a common large consumption endowment is
\[
 W(A)=V_1e_1+V_2e_2-\frac12(e_1^2+e_2^2)-dAe_1.
 \tag{2}
\]

\begin{theorem}[Exact displacement of rival research]
For every $A\in[0,2]$, the research market defined by (1) and $e_1+e_2=1$ has the unique wage and allocation
\[
 w=A/2,\qquad e_1=A/2,\qquad e_2=1-A/2.
 \tag{3}
\]
Consequently an increase in the incumbent prize increases its innovation probability and the research wage while reducing entrant innovation one for one. Total expected successes remain one.
\end{theorem}
\begin{proof}
Strict concavity gives demands $e_1=[A-w]_0^1$ and $e_2=[1-w]_0^1$, where the brackets truncate to $[0,1]$. Substituting $w=A/2$ gives the quantities in (3), both feasible, and their sum is one. For $0<A<2$, both quantities are strictly between zero and one, so the demands have total slope $-2$ near this clearing wage. Since their total is globally nonincreasing, no other wage can clear: any additional root would force demand to remain identically one throughout the interval between them, contradicting the strict local slope. At $A=0$, every $w>0$ produces total demand strictly below one, and $w=0$ clears. At $A=2$, wages below one give total demand above one and wages above one give demand below one. Thus the boundary wages are unique as well. All comparative statements follow immediately from (3).
\end{proof}

\begin{theorem}[Optimal protection in the benchmark]
Under the same assumptions, welfare is a strictly concave quadratic in $A$ and its maximizer on $[0,2]$ is
\[
 A^*=\left[\frac{V_1-V_2+1}{1+2d}\right]_0^2.
 \tag{4}
\]
For any finite nonempty set of feasible durations, the optimal duration minimizes the distance from $A(T)$ to the unconstrained quadratic vertex $(V_1-V_2+1)/(1+2d)$. Ties in distance are genuine welfare ties.
\end{theorem}
\begin{proof}
Substitute (3) into (2). The real research cost is
\[
 \frac12\left((A/2)^2+(1-A/2)^2\right)
 =\frac12-\frac A2+\frac{A^2}{4}.
\]
It follows that
\[
 W(A)=V_2-\frac12+
 \frac{V_1-V_2+1}{2}A-
 \left(\frac14+\frac d2\right)A^2.
\]
The coefficient of $A^2$ is strictly negative. Differentiation and projection onto the interval prove (4). Completing the square shows that maximizing this same expression over a finite set is equivalent to minimizing squared distance to its unconstrained vertex. Finiteness guarantees an attained optimum.
\end{proof}

\paragraph{A policy comparison and its limits.}
If $V_1=V_2=1$ and $d=1/2$, equation (4) gives $A^*=1/2$. Raising the prize from one to two doubles incumbent research from $1/2$ to one, eliminates entrant research, and reduces welfare from $1/2$ to $-1/2$. These welfare figures are changes from the common endowment, so a negative value does not imply negative feasible consumption. If instead $V_1<V_2$, even the quality-weighted output term $V_1e_1+V_2e_2$ decreases with protection despite the incumbent's higher success probability.

The result makes a partial-equilibrium extrapolation fail for an explicit reason: a researcher attracted to one project is unavailable to another. Counting the wage increase as an additional social research cost would double count a transfer and distort (4). Conversely, making research supply elastic would add a real opportunity cost and change (3), so the fixed-supply formula cannot simply be carried into that extension.

\paragraph{Residual gap.}
This research market has no merger choice, strategic acquisition, endogenous entrant prize, or repeated replacement of technology leaders. It does not establish an optimal patent duration for a creative-destruction economy. The finite-duration qualification also matters: an unbounded sequence of patent lengths may approach a limiting prize without attaining it. The solved contribution is an internally accounted equilibrium channel through which stronger rents reallocate scarce research and need not improve welfare; the broader joint competition-policy and intellectual-property problem remains unresolved.

\section{Completion Estimate}
25\%

This is a post hoc, heuristic estimate for the full catalogue target.
The attempt solves research-labor clearing and optimal patent rewards in a two-project market, quantifying crowding out and nonmonotone welfare. Joint patent and competition policy in a dynamic creative-destruction economy still requires endogenous entry, acquisition, diffusion, repeated replacement and identified innovation responses.

\begin{thebibliography}{9}
\bibitem{akcigit} U. Akcigit (2026), ``Creative destruction in economic growth,'' \emph{Scandinavian Journal of Economics}, Section 7, question 3. \href{https://onlinelibrary.wiley.com/doi/10.1111/sjoe.70037}{Publisher article}.
\end{thebibliography}

\end{document}
